\section{Programa completo: \code{sw/game.s}}
\label{anx:game}

\lstinputlisting[language=rvasm]{anexos/game.s}

\newpage
\section{Listado de traducción: \code{sw/game.lst}}
\label{anx:listado}

Generado con \code{python3 asm/pochoasm.py sw/game.s -o sw/game.hex -l sw/game.lst}.
Cada línea tiene la dirección, la palabra de 32 bits que queda en \code{game.hex} y la
línea de \code{game.s} que la produjo. Las pseudo-instrucciones que se expanden a dos
palabras (\code{li} de valores grandes) muestran el texto solo en la primera.

{\scriptsize
\lstinputlisting[language=rvasm, basicstyle=\ttfamily\scriptsize]{anexos/game.lst}
}
